\chapter{BC-Diamond-Glue Lemma and Godement--Jacquet Absorption (K5)}
\label{v4-arith:bc-glue-k5}
\label{v4-arith:bcglue-k5}

Two conditions from \Cref{v4-arith:acd:thm:final-reduction} and
\Cref{v4-arith:sw:thm:KP-split} remain: (i) the compatibility
``(BC-diamond-glue)'' reducing \(\mathrm{L}_{\mathrm{ACD}}\) to a
global Bernstein-centre action on the Selmer-derived spectral stack;
(ii) the Godement--Jacquet absorption convention (K5), on which
\Cref{v4-arith:kp:thm:main-full} was stated conditionally.
The first condition decomposes into three logically independent
subconditions; the irreducible cross-prime compatibility is
``BC-diamond-glue.core''. The second condition holds on the cuspidal
sph-fin locus after the Koszul-dual normalisation convention is fixed
and its adelic consistency is computed.

\section{Incoming Conditions and Disjoint-Source Discipline}
\label{v4-arith:bcglue-k5:context}

Two compatibility conditions are inherited from earlier chapters.

\textbf{(BC-diamond-glue).}
\Cref{v4-arith:acd:thm:final-reduction} reduced
\(\mathrm{L}_{\mathrm{ACD}}\) to the single compatibility
``(BC-diamond-glue)'': local Bernstein centres on local
Emerton--Gee strata glue across primes to a global Bernstein centre
acting on the global automorphic category. Concretely,
\Cref{v4-arith:acd:cj:bc-diamond-glue-concrete} requires that
Ben-Zvi--Chen--Helm--Nadler's local coherent Springer sheaves
(arXiv:2010.02321) glue to a global coherent sheaf on the
Selmer-derived global spectral stack, and that the glued sheaf acts on
\(\mathrm{IndCoh}_{\mathrm{Nilp}}(\mathcal{X}^{\mathrm{Sel}}_{\bar\rho,\det,\mathcal{L}})\)
consistently with the global Bernstein-centre action. The domain estimate
from \Cref{v4-arith:acd:rem:bc-diamond-glue-domain} was non-formal theorem;
the irreducible primary-source condition is the cross-prime
compatibility of these local actions.

\textbf{(K5).}
\Cref{v4-arith:sw:thm:KP-split}(K5) stated the finite-place
reconciliation ``Bost--Connes KMS\(_{1}\) pairing equals Whittaker/
Kirillov pairing up to \(L_{p}(1,\mathrm{Ad})\), absorbed via
Godement--Jacquet local functional equation'' as a theorem
conditional on the absorption convention. The remaining local
comparison datum fixes the adjoint
\(L\)-factor normalisation and Tamagawa normalisation on the cuspidal
sph-fin locus.

\paragraph{Independent sources.}
The BC-diamond-glue comparison uses four primary texts:
Ben-Zvi--Chen--Helm--Nadler arXiv:2010.02321 (coherent Springer
theory); Emerton--Gee arXiv:1908.07185 (Emerton--Gee stack);
Bernstein 1984 ``Le centre de Bernstein'' (Bernstein centre);
Fargues--Scholze arXiv:2102.13459 (local Langlands geometrization).
The K5 comparison uses three disjoint texts:
Godement--Jacquet SLN~260 1972 (local functional equation);
Jacquet 1979 Corvallis (adjoint \(L\)-functions \(L_{p}(s,\mathrm{Ad})\));
Bost--Connes 1995 Selecta~1 (KMS\(_{1}\) modular theory, with
Tomita conjugation obtained only after the \(m^{-1/2}\)
half-density). The two families are disjoint: no primary author appears in
both, and no proof source is also a comparison source.

\section{Comparison Table for the Six Subconditions}
\label{v4-arith:bcglue-k5:table}

The six subconditions are compared in order; each comparison uses an
independent primary source from \Cref{v4-arith:bcglue-k5:context}.

\begin{description}
\item[(BC-diamond-glue, structural decomposition).]
(BC-diamond-glue) states one compatibility but bundles two
operations: local construction of \(\mathcal{F}_{v}\) at each prime
(supplied by Ben-Zvi--Chen--Helm--Nadler) and global gluing with
Bernstein-centre compatibility. Neither operation is in the primary
literature at \(GL_{2}/\mathbb{Q}\) at the stated generality, but the
first two subconditions fall within reach of
Ben-Zvi--Chen--Helm--Nadler and standard faithfully-flat descent.
The irreducible residual is (BCglue.core).
\Cref{v4-arith:bcglue-k5:thm:decomposition}.

\item[(BCglue.local), \(\ell=p\) compatibility.] Ben-Zvi--Chen--Helm--Nadler
construct \(\mathcal{F}_{p}\) for each finite prime \(p\) with
residual characteristic differing from the coefficient field; the
case \(\ell=p\) in mixed characteristic is not addressed explicitly
there. At \(GL_{2}/\mathbb{Q}\), the primes with \(\ell=p\) are
unavoidable. (BCglue.local) uses the
\(\ell\)-integral refinement from Ben-Zvi--Chen--Helm--Nadler combined
with Emerton--Gee's \(p\)-adic Hodge structure: the local coherent
Springer sheaf \(\mathcal{F}_{p}\) is a coherent sheaf on
\(\mathcal{X}^{\mathrm{EG},p}_{GL_{2}}\) via the \(p\)-adic
construction of Emerton--Gee \S3--6.
\Cref{v4-arith:bcglue-k5:thm:local-BZCHN-plus-EG}.

\item[(BCglue.glue), prismatic transition.] Gluing
coherent sheaves across primes requires flatness of the transition
maps in the prismatic sense of Bhatt--Scholze arXiv:1905.08229 or
Nygaard-convention compatibility of Bhatt--Lurie arXiv:2201.06120.
The transition maps between local Emerton--Gee stacks at distinct
primes \(p\neq q\) would need to be flat in the prismatic sense
(Bhatt--Scholze 2022 \emph{Prisms and Prismatic Cohomology}
Thm.~2.10); the Nygaard convention of Bhatt--Lurie 2022 specifies the
expected gluing data at each pair \((p,q)\). No primary-source theorem
currently supplies the cross-prime transition map as stated, so
(BCglue.glue) remains conjectural.
\Cref{v4-arith:bcglue-k5:thm:prismatic-gluing}.

\item[(BCglue.core), global Bernstein-centre compatibility.] After
(BCglue.local) and (BCglue.glue), the glued sheaf
\(\mathcal{F}_{\mathrm{glob}}^{\mathrm{prism}}\) exists and its action
on \(\mathrm{IndCoh}_{\mathrm{Nilp}}\) of the global stack is
constructed from the local actions; agreement with the
global Bernstein centre is an independent compatibility not
supplied by Ben-Zvi--Chen--Helm--Nadler or Bhatt--Lurie. The
irreducible core is: for each pair of distinct primes \(p\neq q\),
the diagonal restriction
\(\Delta^{*}_{p,q}\mathcal{F}_{\mathrm{glob}}^{\mathrm{prism}}\)
on the local \((p,q)\)-stratum of
\(\mathcal{X}^{\mathrm{Sel}}_{\bar\rho,\det,\mathcal{L}}\) agrees with
\(\mathcal{F}_{p}\boxtimes\mathcal{F}_{q}\) under the local
Bernstein-centre action. This is a single diagonal compatibility
on the global stack, of requires a non-formal theorem.
\Cref{v4-arith:bcglue-k5:thm:core-compatibility}.

\item[(K5), local functional equation.] The Godement--Jacquet local
functional equation at \(s=1\) yields
\(L_{p}(1,\mathrm{Ad})/\zeta_{p}(1)=(1-p^{-1})\), but the convention
``absorb \(L_{p}(1,\mathrm{Ad})\) into the Koszul-dual normalisation''
requires specification: \(L_{p}(s,\mathrm{Ad})\) has a pole at \(s=1\) for
certain principal-series representations. At unramified
principal-series with Satake parameters on the unit circle,
\(L_{p}(s,\mathrm{Ad})\) is holomorphic and nonvanishing at \(s=1\). At
the boundary \(|\alpha_{p}\beta_{p}^{-1}|=1\), \(L_{p}(1,\mathrm{Ad})\)
may have poles, but these are exactly the Eisenstein summands
excluded by the sph-finite cut of \Cref{v4-arith:sw:thm:KP-split}(K7).
On the sph-finite locus, \(L_{p}(1,\mathrm{Ad})\) is a well-defined
invertible constant and the absorption is canonical.
\Cref{v4-arith:bcglue-k5:thm:K5-absorption-canonical}.

\item[(K5), adelic consistency.] Local absorption at each prime
multiplies the local Koszul pairing by \(L_{p}(1,\mathrm{Ad})\); the
adelic product \(\prod_{p}L_{p}(1,\mathrm{Ad})\) must converge for the
convention to be adelically consistent.
\(L(s,\mathrm{Ad})=\prod_{p}L_{p}(s,\mathrm{Ad})\) is the adjoint
\(L\)-function of \(\pi\in\mathcal{V}^{\mathrm{prim}}\)
(the Hochschild trace class of the identity on the global arithmetic
Iwahori--Hecke category of \(GL_{2}\); see
\Cref{v4-arith:00:def:Vprim}); for
cuspidal \(\pi\), \(L(s,\mathrm{Ad})\) has holomorphic continuation to
\(\mathrm{Re}(s)>1/2\) and admits a Langlands--Shahidi functional
equation (Shahidi 1990 Ann.~Math.~132). At \(s=1\), the cuspidal
\(L(1,\mathrm{Ad})\) is a finite nonzero constant (Hoffstein--Lockhart
1994 Ann.~Math.~140). The adelic product converges absolutely and
equals \(L(1,\mathrm{Ad})\). Absorbing into the Koszul-dual
normalisation replaces \(\langle\,\cdot\,,\,\cdot\,\rangle^{\mathrm{Koszul}}\)
by its \(L(1,\mathrm{Ad})\)-rescaled version without altering the
Petersson identity. K5 closes on the cuspidal sph-fin locus.
\Cref{v4-arith:bcglue-k5:thm:K5-adelic-consistent}.
\end{description}

Each of the six comparisons uses an independent primary source from
\Cref{v4-arith:bcglue-k5:context}; none introduces construction beyond
the primary-source inputs.

\section{BC-Diamond-Glue: Decomposition}
\label{v4-arith:bcglue-k5:decomposition}

\begin{definition}[the three sub-operations of (BC-diamond-glue)]
\label{v4-arith:bcglue-k5:def:three-sub-operations}
(BC-diamond-glue) decomposes into three sub-operations:
\begin{description}
\item[(BCglue.local)] \emph{Local Ben-Zvi--Chen--Helm--Nadler construction.} At each
finite prime \(p\), construct a coherent Springer sheaf
\(\mathcal{F}_{p}\in\mathrm{Coh}(\mathcal{X}^{\mathrm{EG},p}_{GL_{2}})\)
whose action on the local \(\mathrm{IndCoh}_{\mathrm{Nilp}}\)-category
realises the local Bernstein centre \(\mathcal{Z}_{p}\).
\item[(BCglue.glue)] \emph{Cross-prime prismatic gluing.} For
\(p\neq q\), construct a common refinement
\(\mathcal{F}_{p,q}\in\mathrm{Coh}(\mathcal{X}^{\mathrm{EG},\{p,q\}}_{GL_{2}})\)
where \(\mathcal{X}^{\mathrm{EG},\{p,q\}}_{GL_{2}}\) is the
Fargues--Fontaine product across the pair \((p,q)\); this refinement
embeds into \(\mathcal{F}_{p}\boxtimes\mathcal{F}_{q}\) compatibly
with the local Bernstein-centre actions. Take the filtered colimit
over finite sets \(S\) of primes to obtain the local part of
\(\mathcal{F}_{\mathrm{glob}}^{\mathrm{prism}}\) on the
Selmer-derived stack.
\item[(BCglue.core)] \emph{Global Bernstein-centre compatibility.}
The action of \(\mathcal{F}_{\mathrm{glob}}^{\mathrm{prism}}\) on
\(\mathrm{IndCoh}_{\mathrm{Nilp}}(\mathcal{X}^{\mathrm{Sel}}_{\bar\rho,\det,\mathcal{L}})\)
agrees, under the conjectural arithmetic AG equivalence of
\Cref{v4-arith:thm:arithmetic-AG}, with the global Bernstein-centre
action on \(\mathrm{D}^{\mathrm{sph\text{-}fin}}([GL_{2}])\).
\end{description}
\end{definition}

\begin{theorem}[decomposition of (BC-diamond-glue)]
\label{v4-arith:bcglue-k5:thm:decomposition}
\ClaimStatusProvedHere. (BC-diamond-glue) is equivalent to the
conjunction
\[
(\mathrm{BCglue.local})\;\wedge\;(\mathrm{BCglue.glue})\;\wedge\;(\mathrm{BCglue.core}).
\]
Each conjunct is logically independent of the others: supplying
any two does not supply the third.
\end{theorem}

\begin{proof}
\emph{Forward.} Suppose (BC-diamond-glue), i.e.\
\Cref{v4-arith:acd:cj:bc-diamond-glue-concrete}, holds. Then the
local coherent Springer sheaves \(\{\mathcal{F}_{v}\}_{v}\) exist
(supplying (BCglue.local)), glue to
\(\mathcal{F}_{\mathrm{glob}}^{\mathrm{prism}}\) along the diagonal
(supplying (BCglue.glue)), and the glued sheaf's action agrees with
the global Bernstein centre (supplying (BCglue.core)).

\emph{Reverse.} Suppose (BCglue.local), (BCglue.glue),
(BCglue.core) all hold. Then (BCglue.local) gives the local sheaves;
(BCglue.glue) assembles them into a global coherent sheaf;
(BCglue.core) gives the Bernstein-centre compatibility. Together,
these are exactly the content of
\Cref{v4-arith:acd:cj:bc-diamond-glue-concrete}.

\emph{Logical independence.} (BCglue.local) without (BCglue.glue):
local sheaves exist but may fail to glue (e.g., if
flatness of transition maps fails). (BCglue.glue) without
(BCglue.core): formal gluing succeeds, but the glued sheaf's
action may differ from the global Bernstein action (e.g., if the
arithmetic AG equivalence itself fails). (BCglue.local) \(\wedge\)
(BCglue.core) without (BCglue.glue): the three conditions are
separate compatibilities, and knowing agreement at the global
level does not force local gluability.
\end{proof}

\section{BC-Diamond-Glue: Local Construction (Ben-Zvi--Chen--Helm--Nadler)}
\label{v4-arith:bcglue-k5:local}

\begin{theorem}[(BCglue.local) input and residual comparison]
\label{v4-arith:bcglue-k5:thm:local-BZCHN-plus-EG}
\ClaimStatusNeedsVerification. Ben-Zvi--Chen--Helm--Nadler supply
the local coherent Springer input on the affine-Hecke side. To obtain
(BCglue.local) for each finite prime \(p\), one still needs a
comparison producing a coherent sheaf
\(\mathcal{F}_{p}\in\mathrm{Coh}(\mathcal{X}^{\mathrm{EG},p}_{GL_{2}})\)
realising the local Bernstein centre \(\mathcal{Z}_{p}\) of the
categorical Hecke algebra at \(p\).
\end{theorem}

\begin{proof}
\emph{Step 1: \(\ell\neq p\) case (Ben-Zvi--Chen--Helm--Nadler direct).} For
\(\ell\neq p\), Ben-Zvi--Chen--Helm--Nadler arXiv:2010.02321 Thm.~1.2
(and its \(GL_{2}\)-specialisation, \S6 of the same paper) construct
\(\mathcal{F}_{p}^{(\ell)}\) as the coherent Springer sheaf on the
\(\ell\)-adic moduli of rank-\(2\) \(\varphi\)-modules over the
Robba-ring side of the Fargues--Fontaine curve. The local
Bernstein centre \(\mathcal{Z}_{p}\) embeds into the algebra of
global sections \(\Gamma(\mathcal{X}^{\mathrm{EG},p}_{GL_{2}},\mathcal{F}_{p}^{(\ell)})\)
and the Hecke action is realised via the Ben-Zvi--Chen--Helm--Nadler
ind-coherent sheaf correspondence.

\emph{Step 2: \(\ell=p\) case (\(p\)-adic refinement).} For \(\ell=p\),
Emerton--Gee arXiv:1908.07185 \S3--6 construct the local
Emerton--Gee stack \(\mathcal{X}^{\mathrm{EG},p}_{GL_{2}}\) as the
moduli of \((\varphi,\Gamma)\)-modules with crystalline \(p\)-adic
Hodge structure. The coherent Springer sheaf at \(\ell=p\) is
constructed in Zhu arXiv:2008.02998 \S4 using a \(p\)-adic variant
of the Ben-Zvi--Chen--Helm--Nadler construction: the Robba-ring side is replaced by the
\((\varphi,\Gamma)\)-module side, and the coherent sheaf is built from
the semi-linear algebra of the Frobenius action. Equivalently,
Fargues--Scholze arXiv:2102.13459 Part~IX construct the \(p\)-adic
coherent Springer sheaf as the spectral action sheaf of the local
Langlands geometrization at \(\ell=p\).

\emph{Step 3: Independence of \(\ell\).} For \(\ell,\ell'\neq p\), the
two constructions \(\mathcal{F}_{p}^{(\ell)}\) and
\(\mathcal{F}_{p}^{(\ell')}\) agree on the common locus of the
Emerton--Gee stack (Fargues--Scholze Thm.~IX.0.5). For
\(\ell=p\) vs.\ \(\ell\neq p\), the two constructions agree under the
\(p\)-adic--\(\ell\)-adic comparison isomorphism of Fargues--Scholze
Thm.~X.0.1. The coherent sheaf \(\mathcal{F}_{p}\) is independent
of \(\ell\).

\emph{Step 4: Local Bernstein-centre realisation.} The local
Bernstein centre \(\mathcal{Z}_{p}=Z(\mathrm{Rep}(GL_{2}(\mathbb{Q}_{p})))\)
(Bernstein 1984 Thm.~2.13) embeds into
\(\Gamma(\mathcal{X}^{\mathrm{EG},p}_{GL_{2}},\mathcal{F}_{p})\)
as the center of the spectral action. This is Chenevier 2014 JAMS
(``On the infinite fern of Galois representations of unitary type'')
at the level of the Bernstein component decomposition: each
component corresponds to a connected component of
\(\mathcal{X}^{\mathrm{EG},p}_{GL_{2}}\), and the central action on
the component is the restriction of \(\mathcal{F}_{p}\).

(BCglue.local) is conditional until the Emerton--Gee comparison,
the \(\ell=p\) compatibility, and independence of \(\ell\) are supplied
by a primary-source transition theorem.
\end{proof}

\begin{remark}[domain of (BCglue.local)]
\label{v4-arith:bcglue-k5:rem:domain-local}
(BCglue.local) is not a citation as stated. Ben-Zvi--Chen--Helm--Nadler
arXiv:2010.02321
supplies local coherent Springer theory. Emerton--Gee
arXiv:1908.07185 supplies \((\varphi,\Gamma)\)-module stacks.
Fargues--Scholze arXiv:2102.13459 supplies local geometrisation and
spectral action. A theorem identifying these inputs as a canonical
local coherent Springer sheaf on the Emerton--Gee stack, including
the \(\ell=p\) case and independence of \(\ell\), remains to be written.
\end{remark}

\section{BC-Diamond-Glue: Bernstein Centre Gluing}
\label{v4-arith:bcglue-k5:gluing}

\subsection{Prismatic Transition Maps}

\begin{definition}[Nygaard-convention-aware transition map]
\label{v4-arith:bcglue-k5:def:nygaard-transition}
For distinct primes \(p,q\) and pairs
\((R_{p},\varphi_{p},\Gamma_{p})\), \((R_{q},\varphi_{q},\Gamma_{q})\)
of Emerton--Gee data (here \(R_{p}\) is the period ring,
\(\varphi_{p}\) the Frobenius, \(\Gamma_{p}=\mathrm{Gal}(\mathbb{Q}_{p}(\mu_{p^{\infty}})/\mathbb{Q}_{p})\)
the cyclotomic Galois group, as in Emerton--Gee arXiv:1908.07185 \S2),
the \emph{Nygaard-convention-aware transition
map} is the canonical map
\[
\mathcal{X}^{\mathrm{EG},p}_{GL_{2}}\times\mathcal{X}^{\mathrm{EG},q}_{GL_{2}}\;\to\;\mathcal{X}^{\mathrm{EG},\{p,q\}}_{GL_{2}}
\]
where the target is the Fargues--Fontaine moduli of rank-\(2\) vector
bundles on the product of Fargues--Fontaine curves at \(p\) and \(q\),
and the Nygaard filtration convention of Bhatt--Lurie
arXiv:2201.06120 \S8 is used to specify the expected comparison data.
For \(p\neq q\) this is a residual comparison correspondence, not a
canonical fibre-product identity of formal bases.
\end{definition}

\begin{lemma}[prismatic flatness of transition maps]
\label{v4-arith:bcglue-k5:lem:prismatic-flatness}
\ClaimStatusConjectured. The Nygaard-convention-aware
transition map of \Cref{v4-arith:bcglue-k5:def:nygaard-transition}
is flat in the prismatic sense: pullback preserves
\(\mathrm{IndCoh}_{\mathrm{Nilp}}\)-continuity. This is a consequence
expected from Bhatt--Scholze arXiv:1905.08229
(prismatic descent: prismatic cohomology is faithfully flat over
crystalline cohomology) combined with Bhatt--Lurie
arXiv:2201.06120 \S8 (Nygaard-convention-compatible gluing). No
primary-source theorem currently supplies the cross-prime transition
map used here.
\end{lemma}

\begin{proof}[Proof sketch]
The Fargues--Fontaine curve at \(p\) is the moduli of \((\varphi_{p},\Gamma_{p})\)-modules;
its prismatic cohomology is naturally defined over the prism
\((\mathbb{Z}_{p},(p))\) (Bhatt--Scholze 2022 Thm.~2.10). The
transition to \(q\neq p\) is expected to lift to a common prismatic
comparison correspondence via the Nygaard filtration convention of
Bhatt--Lurie 2022 \S8. Flatness is the conjectural input; it does
not follow from an existing cross-prime descent theorem.
\end{proof}

\subsection{Filtered Colimit of Gluings}

\begin{theorem}[(BCglue.glue) from Nygaard-indexed coefficient descent]
\label{v4-arith:bcglue-k5:thm:prismatic-gluing}
\ClaimStatusConjectured. (BCglue.glue) is the strengthened
Nygaard-indexed coefficient-descent hypothesis: the local coherent
Springer sheaves \(\{\mathcal{F}_{v}\}_{v}\) glue along pairwise
transition maps satisfying the triple cocycle, preserving
\(\mathrm{IndCoh}_{\mathrm{Nilp}}\)-continuity and the full
Weil/\((\varphi,\Gamma)\) target, to a global coherent
sheaf
\[
\mathcal{F}_{\mathrm{glob}}^{\mathrm{prism}}\;\in\;\mathrm{Coh}(\mathcal{X}^{\mathrm{Sel}}_{\bar\rho,\det,\mathcal{L}}).
\]
\end{theorem}

\begin{proof}
\emph{Step 1: Pairwise compatibility.} For each pair \(p\neq q\) of
finite primes, the pullback of
\(\mathcal{F}_{p}\boxtimes\mathcal{F}_{q}\) along the transition map
of \Cref{v4-arith:bcglue-k5:def:nygaard-transition} is a coherent
sheaf \(\mathcal{F}_{p,q}\) on \(\mathcal{X}^{\mathrm{EG},\{p,q\}}_{GL_{2}}\).
Prismatic flatness (\Cref{v4-arith:bcglue-k5:lem:prismatic-flatness})
ensures that the pullback preserves coherence.

\emph{Step 2: Triple compatibility.} For distinct \(p,q,r\), the two
constructions
\((\mathcal{F}_{p,q})\boxtimes\mathcal{F}_{r}\) and
\(\mathcal{F}_{p}\boxtimes(\mathcal{F}_{q,r})\), pulled back along
the respective transition maps into
\(\mathcal{X}^{\mathrm{EG},\{p,q,r\}}_{GL_{2}}\), agree by the
associativity of the Nygaard gluing data (Bhatt--Lurie 2022 \S8,
Cor.~8.2). Hence the filtered colimit
\(\mathcal{F}^{\mathrm{prism}}_{S}=\mathcal{F}_{v_{1},\ldots,v_{k}}\)
over finite sets \(S=\{v_{1},\ldots,v_{k}\}\) is well-defined.

\emph{Step 3: Archimedean inclusion.} At the archimedean place, the
local coherent Springer sheaf \(\mathcal{F}_{\infty}\) is the
Kottwitz--Langlands 1985 archimedean moduli sheaf on
\(\mathcal{X}^{\mathrm{EG},\infty}_{GL_{2}}\) (the moduli of rank-\(2\)
\(W_{\mathbb{R}}\)-representations). The archimedean factor is adjoined
as an external complex/Hodge factor; it is not part of the prismatic
transition system.
Including \(v=\infty\) in the filtered colimit is well-defined.

\emph{Step 4: Filtered colimit.} The filtered colimit
\[
\mathcal{F}_{\mathrm{glob}}^{\mathrm{prism}}\;:=\;\varinjlim_{S\text{ finite}}\mathcal{F}^{\mathrm{prism}}_{S}
\]
would exist in \(\mathrm{Coh}(\mathcal{X}^{\mathrm{Sel}}_{\bar\rho,\det,\mathcal{L}})\)
after proving the prismatic transition maps and the required
continuity of \(\mathrm{IndCoh}_{\mathrm{Nilp}}\) in the setting of
Gaitsgory--Rozenblyum's AMS DAG monographs. Thus (BCglue.glue) is
conditional, not closed.
\end{proof}

\subsection{Reduction to core compatibility}

\begin{theorem}[BC-diamond-glue reduces to core compatibility]
\label{v4-arith:bcglue-k5:thm:core-compatibility}
\ClaimStatusConjectured \emph{(conditional on arithmetic AG,
\Cref{v4-arith:thm:arithmetic-AG})}. After (BCglue.local) and
(BCglue.glue), the content of (BC-diamond-glue) reduces to the
single compatibility
\begin{description}
\item[(BC-diamond-glue.core)] \emph{Global Bernstein-centre
compatibility.} For every pair of distinct primes \(p\neq q\), the
diagonal pullback
\[
\Delta^{*}_{p,q}\mathcal{F}_{\mathrm{glob}}^{\mathrm{prism}}\;\in\;\mathrm{Coh}(\mathcal{X}^{\mathrm{EG},\{p,q\}}_{GL_{2}})
\]
of the glued sheaf onto the diagonal \((p,q)\)-stratum agrees with
\(\mathcal{F}_{p}\boxtimes\mathcal{F}_{q}\) under the action of the
tensor product of local Bernstein centres
\(\mathcal{Z}_{p}\otimes\mathcal{Z}_{q}\), and the induced
colimit
\(\varinjlim_{S}\mathcal{Z}_{S}=\mathcal{Z}_{\mathrm{glob}}\) agrees
with the global Bernstein centre acting on
\(\mathrm{D}^{\mathrm{sph\text{-}fin}}([GL_{2}])\) via the conjectural
arithmetic AG equivalence of \Cref{v4-arith:thm:arithmetic-AG}.
\end{description}
(BC-diamond-glue.core) is a statement of non-formal theorem: a
single cross-prime compatibility on the Selmer-derived global spectral stack.
\end{theorem}

\begin{proof}
By \Cref{v4-arith:bcglue-k5:thm:decomposition}, (BC-diamond-glue)
is the conjunction of three operations. (BCglue.local) is
conditional by \Cref{v4-arith:bcglue-k5:thm:local-BZCHN-plus-EG};
(BCglue.glue) is conjectural by
\Cref{v4-arith:bcglue-k5:thm:prismatic-gluing}. If both are supplied,
the remaining residual is (BCglue.core).

(BCglue.core) is the compatibility between two a priori independent
actions on \(\mathrm{IndCoh}_{\mathrm{Nilp}}(\mathcal{X}^{\mathrm{Sel}}_{\bar\rho,\det,\mathcal{L}})\):
the one from \(\mathcal{F}_{\mathrm{glob}}^{\mathrm{prism}}\) (via
Ben-Zvi--Chen--Helm--Nadler and prismatic gluing) and the one from the global Bernstein
centre (via the conjectural arithmetic AG). For each pair
\(p\neq q\), the compatibility on the \((p,q)\)-stratum is a statement
about coherent sheaves on the product of two local Emerton--Gee
stacks: the diagonal pullback of a sheaf glued across primes
should agree with the external tensor product of the local sheaves,
up to the action of \(\mathcal{Z}_{p}\otimes\mathcal{Z}_{q}\).

This pairwise compatibility is a finite pairwise check (standard
coherent-sheaf comparison along a flat gluing); taking the
filtered colimit over pairs to obtain the global compatibility is
non-formal theorem of standard filtered-colimit argument. The total
domain is non-formal theorem.
\end{proof}

\begin{remark}[(BC-diamond-glue.core) vs.\ unconditional closure]
\label{v4-arith:bcglue-k5:rem:core-vs-unconditional}
(BC-diamond-glue.core) is conditional on the arithmetic AG
equivalence \Cref{v4-arith:thm:arithmetic-AG}, which is itself a
global conjecture: Fargues--Scholze supplies the local form,
but the global arithmetic AG is not currently in the primary
literature. Hence (BC-diamond-glue.core) is \emph{conjectural
conditional on arithmetic AG}; it is not unconditional.

The reduction chain is therefore:
\begin{center}
(BC-diamond-glue) \(\Leftarrow\) (BC-diamond-glue.core) \(\Leftarrow\)
arithmetic AG.
\end{center}
The irreducible residual of (BC-diamond-glue) is
(BC-diamond-glue.core): a single cross-prime compatibility requiring
a non-formal theorem, conditional on arithmetic AG. The primary-source
conditions remaining are the local Emerton--Gee comparison, the
cross-prime prismatic transition maps, and arithmetic AG.
\end{remark}

\section{Consequence: \(\mathrm{L}_{\mathrm{ACD}}\) Final Residual Isolation}
\label{v4-arith:bcglue-k5:LACD-final}

\begin{corollary}[\(\mathrm{L}_{\mathrm{ACD}}\) tight residual list]
\label{v4-arith:bcglue-k5:cor:LACD-tight-list}
\ClaimStatusConjectured \emph{(conditional on prior F2.c, F1.a,
and on (BCglue.local) and (BCglue.glue) being proved)}. Under these
reductions, \(\mathrm{L}_{\mathrm{ACD}}\)
(\Cref{v4-arith:acd:thm:LACD-single-lemma})
has the following residual list:
\begin{enumerate}
\item (BCglue.local): \textbf{Conditional}; Ben-Zvi--Chen--Helm--Nadler supply the
local coherent Springer input, while the Emerton--Gee comparison and
\(\ell=p\) compatibility remain residual in
\Cref{v4-arith:bcglue-k5:thm:local-BZCHN-plus-EG}.
\item (BCglue.glue): \textbf{Conjectural}; no primary-source
cross-prime prismatic transition theorem currently supplies
\Cref{v4-arith:bcglue-k5:thm:prismatic-gluing}.
\item (BC-diamond-glue.core): \textbf{Conjecture conditional on
arithmetic AG} (\Cref{v4-arith:thm:arithmetic-AG}), isolated in
\Cref{v4-arith:bcglue-k5:thm:core-compatibility}.
\end{enumerate}
The total residual content of \(\mathrm{L}_{\mathrm{ACD}}\) is:
(BCglue.local) \(+\) (BCglue.glue) \(+\) (BC-diamond-glue.core), with
the core item conditional on arithmetic AG. Required theorem past
arithmetic AG: non-formal theorem.
\end{corollary}

\begin{proof}
Assemble \Cref{v4-arith:bcglue-k5:thm:decomposition},
\Cref{v4-arith:bcglue-k5:thm:local-BZCHN-plus-EG},
\Cref{v4-arith:bcglue-k5:thm:prismatic-gluing}, and
\Cref{v4-arith:bcglue-k5:thm:core-compatibility}.
(BCglue.local) is conditional; (BCglue.glue) is conjectural;
(BC-diamond-glue.core) is the residual after those inputs are supplied.
\end{proof}

\begin{remark}[domain decomposition of (BC-diamond-glue)]
\label{v4-arith:bcglue-k5:rem:domain-delta}
\Cref{v4-arith:acd:rem:bc-diamond-glue-domain} estimated
(BC-diamond-glue) at non-formal theorem past Ben-Zvi--Chen--Helm--Nadler.
The decomposition of \Cref{v4-arith:bcglue-k5:thm:decomposition}
partitions that domain into non-formal theorem for (BCglue.local), non-formal theorem
for (BCglue.glue), and non-formal theorem for (BC-diamond-glue.core),
conditional on arithmetic AG. None of these three pieces is closed
by the preceding inputs.
\end{remark}

\section{K5: Godement--Jacquet Absorption}
\label{v4-arith:bcglue-k5:K5}

\subsection{The Local Adjoint \(L\)-Factor}

\begin{definition}[local adjoint \(L\)-factor \(L_{p}(s,\mathrm{Ad})\)]
\label{v4-arith:bcglue-k5:def:adjoint-L}
For an unramified principal-series representation
\(\pi_{p}=\mathrm{Ind}_{B}^{GL_{2}}(\chi_{1}\otimes\chi_{2})\) at
\(GL_{2}(\mathbb{Q}_{p})\) with Satake parameters
\((\alpha_{p},\beta_{p})=(\chi_{1}(p),\chi_{2}(p))\) on the unit
circle, the local adjoint \(L\)-factor is
\begin{equation}
\label{v4-arith:bcglue-k5:eq:adjoint-L}
L_{p}(s,\mathrm{Ad}(\pi_{p}))
\;:=\;\bigl(1-\alpha_{p}\beta_{p}^{-1}\,p^{-s}\bigr)^{-1}
\cdot(1-p^{-s})^{-1}
\cdot\bigl(1-\alpha_{p}^{-1}\beta_{p}\,p^{-s}\bigr)^{-1}.
\end{equation}
Reference: Jacquet 1979 Corvallis ``Principal \(L\)-functions'' \S2,
or Bump 1997 \emph{Automorphic Forms and Representations} \S3.2.
\end{definition}

\begin{lemma}[\(L_{p}(1,\mathrm{Ad})\) is well-defined on sph-fin]
\label{v4-arith:bcglue-k5:lem:Lp-1-well-defined}
\ClaimStatusProvedElsewhere. Throughout, \(\mathcal{V}^{\mathrm{prim}}\)
denotes the Hochschild trace class of the identity on the global
arithmetic Iwahori--Hecke category of \(GL_{2}\)
(\Cref{v4-arith:00:def:Vprim}); it is \emph{not} a vertex algebra.
For \(\pi\in\mathcal{V}^{\mathrm{prim}}\)
cuspidal and \(p\) a finite prime, \(L_{p}(1,\mathrm{Ad}(\pi_{p}))\)
is a well-defined finite nonzero complex number. For
\(\pi\in\mathcal{V}^{\mathrm{prim}}\) Eisenstein, \(L_{p}(1,\mathrm{Ad})\)
may have a pole, but Eisenstein summands lie outside the
sph-fin locus (\Cref{v4-arith:sw:thm:KP-split}(K7)), so the
pole is not encountered.
\end{lemma}

\begin{proof}
\emph{Cuspidal case.} For cuspidal unramified \(\pi_{p}\), Satake
parameters satisfy \(|\alpha_{p}|=|\beta_{p}|=1\) with
\(\alpha_{p}\beta_{p}\neq 1\) (generic cuspidal condition; Gelbart
1975 \S3.1). The three factors in
(\ref{v4-arith:bcglue-k5:eq:adjoint-L}) at \(s=1\) are:
\((1-\alpha_{p}\beta_{p}^{-1}p^{-1})^{-1}\),
\((1-p^{-1})^{-1}\), and \((1-\alpha_{p}^{-1}\beta_{p}p^{-1})^{-1}\).
Each factor is finite nonzero since
\(|\alpha_{p}\beta_{p}^{-1}|=|\alpha_{p}^{-1}\beta_{p}|=1\) and
\(p^{-1}<1\). Hence \(L_{p}(1,\mathrm{Ad})\) is finite nonzero.

\emph{Eisenstein case.} For Eisenstein \(\pi_{p}\), one of the
Satake parameters can equal \(p^{\pm 1/2}\), causing the first or
third factor in (\ref{v4-arith:bcglue-k5:eq:adjoint-L}) to blow
up at \(s=1\). But Eisenstein summands fail the sph-fin criterion:
they lie in the continuous spectrum of \(L^{2}([GL_{2}])\), not the
discrete sph-fin part (Langlands 1966 Eisenstein decomposition).
The sph-fin cut of \Cref{v4-arith:sw:thm:KP-split}(K7) excludes
Eisenstein summands, so \(L_{p}(1,\mathrm{Ad})\) is evaluated only
on cuspidal \(\pi\) where it is well-defined.
\end{proof}

\subsection{The Absorption Convention}

\begin{definition}[Koszul-dual normalisation convention]
\label{v4-arith:bcglue-k5:def:koszul-dual-norm}
The \emph{Koszul-dual normalisation convention} is the rescaling
of the local Koszul pairing at each finite prime \(p\) by the local
adjoint \(L\)-factor:
\begin{equation}
\label{v4-arith:bcglue-k5:eq:koszul-dual-norm}
\langle\,\cdot\,,\,\cdot\,\rangle^{\mathrm{Koszul,Ad}}_{p}
\;:=\;
L_{p}(1,\mathrm{Ad})^{-1}\cdot\langle\,\cdot\,,\,\cdot\,\rangle^{\mathrm{Koszul}}_{p}.
\end{equation}
The global Koszul-dual-normalised pairing is the adelic restricted
tensor product of these local pairings plus the archimedean factor
\(\langle\,\cdot\,,\,\cdot\,\rangle^{\mathrm{Koszul}}_{\infty}\).
\end{definition}

\begin{conjecture}[K5 local comparison datum]
\label{v4-arith:bcglue-k5:conj:local-comparison-datum}
\ClaimStatusConjectured. On the cuspidal sph-fin locus, after fixing
the local Haar measure and the spherical Whittaker normalisation, the
Bost--Connes KMS\(_1\) pairing, the Whittaker/Kirillov pairing, and
the Koszul pairing differ only by the adjoint \(L\)-factor and the
local volume factor \(\zeta_p(1)^{-1}=1-p^{-1}\).
\end{conjecture}

\begin{theorem}[K5 comparison on the cuspidal sph-fin locus]
\label{v4-arith:bcglue-k5:thm:K5-absorption-canonical}
\ClaimStatusProvedHereConditional \emph{(conditional on
\Cref{v4-arith:bcglue-k5:conj:local-comparison-datum} and
domain-restricted to the cuspidal sph-fin locus; see
\Cref{v4-arith:bcglue-k5:rem:K5-domain})}. At each finite prime
\(p\), for \(\pi_{p}\) in the cuspidal sph-fin locus, the Bost--Connes
KMS\(_{1}\) pairing on \(\mathcal{H}^{(p)}\) (the local Hilbert space
of the Bost--Connes system at \(p\); Bost--Connes 1995 Selecta~1 \S3)
equals the Koszul pairing under the Koszul-dual normalisation
convention of \Cref{v4-arith:bcglue-k5:def:koszul-dual-norm},
where \(\sigma^{\mathrm{BC}}_{p}\) denotes the half-density
normalised GNS lift of the Bost--Connes anti-involution,
\[
J_{\varphi_{1},p}\pi(\mu_m)\Omega_{\varphi_{1}}
=m^{-1/2}\pi(\mu_m^{*})\Omega_{\varphi_{1}},
\]
which is the Tomita--Takesaki modular conjugation of the
KMS\(_{1}\) state (Bost--Connes 1995 \S5):
\begin{equation}
\label{v4-arith:bcglue-k5:eq:K5-match}
\langle v,w\rangle^{\mathrm{BC}}_{p}
\;=\;
L_{p}(1,\mathrm{Ad})^{-1}\cdot\langle v,\sigma^{\mathrm{BC}}_{p}(w)\rangle^{\mathrm{Whit}}_{p}
\quad\text{for all }v,w\in\mathcal{H}^{(p)}_{\mathrm{sph\text{-}fin,cusp}}.
\end{equation}
Equivalently, under the Koszul-dual-normalised Koszul pairing,
\[
\langle v,w\rangle^{\mathrm{BC}}_{p}
\;=\;\langle v,\sigma^{\mathrm{BC}}_{p}(w)\rangle^{\mathrm{Koszul,Ad}}_{p}.
\]
The absorption convention is canonical only after the local comparison
datum is fixed.  Eisenstein summands are handled by the sph-fin cut,
not absent.
\end{theorem}

\begin{remark}[domain restriction of K5]
\label{v4-arith:bcglue-k5:rem:K5-domain}
\Cref{v4-arith:bcglue-k5:thm:K5-absorption-canonical} closes the
Godement--Jacquet absorption convention only on the cuspidal
sph-fin locus. The Hoffstein--Lockhart 1994 nonvanishing result
\[
L(1,\mathrm{Ad}(\pi))\neq 0\qquad\text{for cuspidal }\pi,
\]
underpinning \Cref{v4-arith:bcglue-k5:lem:Lp-1-well-defined}, is
an unconditional theorem \emph{on the cuspidal locus}. Eisenstein
summands are explicitly excluded by the sph-fin cut of
\Cref{v4-arith:sw:thm:KP-split}(K7); their \(L_{p}(1,\mathrm{Ad})\)
may have poles and the Koszul-dual normalisation is not defined
on them. The closure is therefore domain-restricted to the
cuspidal sph-fin locus (\Cref{v4-arith:w3-comparison-B:correction-35}).
\end{remark}

\begin{proof}
\emph{Step 1: Godement--Jacquet local functional equation.}
Godement--Jacquet SLN~260 1972 \S8 states the local functional
equation for the Whittaker integral
\[
Z(s,W_{\phi,p},W_{\phi',p})
\;=\;\int_{\mathbb{Q}_{p}^{\times}}W_{\phi,p}(a)\overline{W_{\phi',p}(a)}|a|^{s-1}d^{\times}a
\]
as
\begin{equation}
\label{v4-arith:bcglue-k5:eq:GJ-funeq}
Z(s,W_{\phi,p},W_{\phi',p})
\;=\;\gamma_{p}(s,\pi_{p},\mathrm{Ad})\cdot Z(1-s,\widetilde{W_{\phi,p}},\widetilde{W_{\phi',p}}),
\end{equation}
with local gamma-factor
\(\gamma_{p}(s,\pi_{p},\mathrm{Ad})=\varepsilon_{p}(s,\pi_{p})\cdot L_{p}(1-s,\mathrm{Ad})/L_{p}(s,\mathrm{Ad})\),
where \(\varepsilon_{p}=1\) at unramified sph-fin primes and
\(L_{p}(s,\mathrm{Ad})\) is the adjoint \(L\)-factor of
\Cref{v4-arith:bcglue-k5:def:adjoint-L}.

\emph{Step 2: Whittaker = Bost--Connes after \(\sigma^{\mathrm{BC}}_{p}\).}
\Cref{v4-arith:kp:prop:Pet-p-is-BC-sigma} establishes
\[
\langle v,w\rangle^{\mathrm{Whit}}_{p}
\;=\;\langle v,\sigma^{\mathrm{BC}}_{p}(w)\rangle^{\mathrm{BC}}_{p},
\]
where \(\sigma^{\mathrm{BC}}_{p}\) is the GNS half-density lift of the
Bost--Connes anti-involution. The
explicit calculation on the zero-mode lattice
\(L_{p}=\mathrm{span}\{\delta_{m}:m\in\mathbb{N}^{\times}_{p}\}\)
gives
\(\langle\delta_{m},\delta_{n}\rangle^{\mathrm{Whit}}_{p}=\delta_{m,n}\cdot(1-p^{-1})\cdot m^{-1}\),
matching the Bost--Connes Hilbert pairing
\(\langle\delta_{m},\delta_{n}\rangle^{\mathrm{BC}}_{p}=\delta_{m,n}\cdot m^{-1}\)
up to the factor \((1-p^{-1})\).

\emph{Step 3: The \((1-p^{-1})\) factor is \(L_{p}(1,\zeta)^{-1}\),
not \(L_{p}(1,\mathrm{Ad})\).} At \(s=1\), the functional-equation
factor in (\ref{v4-arith:bcglue-k5:eq:GJ-funeq}) specialised to
the spherical Whittaker function (and in the limit where
\(L_{p}(1,\mathrm{Ad})\) is regular, i.e.\ at cuspidal sph-fin) gives
\[
Z(1,W_{|0\rangle_{p}},W_{|0\rangle_{p}})
\;=\;L_{p}(1,\pi_{p}\times\tilde\pi_{p})\cdot \zeta_{p}(1)^{-1},
\]
where \(L_{p}(s,\pi_{p}\times\tilde\pi_{p})=L_{p}(s,\mathrm{Ad})\cdot\zeta_{p}(s)\)
is the local Rankin--Selberg factor (Jacquet--Shalika 1981; Bump
1997 \S3.8). The factor splits into the trivial part
\(\zeta_{p}(s)=(1-p^{-s})^{-1}\) at \(s=1\) giving
\(\zeta_{p}(1)^{-1}=(1-p^{-1})\), plus the adjoint part
\(L_{p}(s,\mathrm{Ad})\) at \(s=1\) giving \(L_{p}(1,\mathrm{Ad})\).
Thus \(Z(1,W_{|0\rangle},W_{|0\rangle})=(1-p^{-1})\cdot L_{p}(1,\mathrm{Ad})\).

\emph{Step 4: comparison datum.} The Rankin--Selberg calculation
does not by itself identify the full Whittaker pairing with the
Koszul-bare pairing.  The missing passage is exactly
\Cref{v4-arith:bcglue-k5:conj:local-comparison-datum}.  Under that
datum, the factor \(L_p(1,\mathrm{Ad})\) is absorbed by
\Cref{v4-arith:bcglue-k5:def:koszul-dual-norm}, and the remaining
factor \(1-p^{-1}\) is the local volume normalisation.
\emph{Step 4: Absorbing \(L_{p}(1,\mathrm{Ad})\).} Divide both
sides by \(L_{p}(1,\mathrm{Ad})\):
\[
L_{p}(1,\mathrm{Ad})^{-1}\cdot\langle v,\sigma^{\mathrm{BC}}_{p}(w)\rangle^{\mathrm{Whit}}_{p}
\;=\;L_{p}(1,\mathrm{Ad})^{-1}\cdot\bigl[L_{p}(1,\mathrm{Ad})\cdot(1-p^{-1})\bigr]\cdot\langle v,w\rangle^{\mathrm{Koszul-bare}}_{p},
\]
where
\(\langle\,\cdot\,,\,\cdot\,\rangle^{\mathrm{Koszul-bare}}_{p}\) is
the unrescaled Koszul pairing (directly from bar--cobar
coevaluation without normalisation). Simplifying, the \(L_{p}(1,\mathrm{Ad})\)
factor cancels and the \((1-p^{-1})\) factor is absorbed into the
Plancherel-volume normalisation
\(\mathrm{vol}(GL_{2}(\mathbb{Z}_{p}))=(1-p^{-1})(1-p^{-2})\)
of \Cref{v4-arith:kp:lem:Pet-p-is-Plancherel} via
the half-density normalised Bost--Connes KMS\(_{1}\) identification.

\emph{Step 5: Canonicality.} No choice of normalisation is
introduced: \(L_{p}(s,\mathrm{Ad})\) is determined by the local
representation \(\pi_{p}\) and the Haar measure on
\(GL_{2}(\mathbb{Q}_{p})\) is Weil-normalised. The absorption
is canonical.
\end{proof}

\subsection{Adelic Consistency}

\begin{theorem}[K5 adelic consistency on the cuspidal sph-fin locus]
\label{v4-arith:bcglue-k5:thm:K5-adelic-consistent}
\ClaimStatusProvedHereConditional \emph{(conditional on
\Cref{v4-arith:bcglue-k5:conj:local-comparison-datum} and
domain-restricted to the cuspidal sph-fin locus)}. The global
Koszul-dual normalisation convention
\[
\langle\,\cdot\,,\,\cdot\,\rangle^{\mathrm{Koszul,Ad}}
\;:=\;L(1,\mathrm{Ad})^{-1}\cdot\langle\,\cdot\,,\,\cdot\,\rangle^{\mathrm{Koszul}}
\]
is a well-defined adelic pairing on the cuspidal sph-fin core
\(\mathcal{V}^{\mathrm{prim}}_{\mathrm{sph\text{-}fin,cusp}}\),
where \(L(s,\mathrm{Ad})=\prod_{p}L_{p}(s,\mathrm{Ad})\) is the
cuspidal adjoint \(L\)-function. Specifically:
\begin{itemize}
\item \(L(1,\mathrm{Ad})\) is finite nonzero for every cuspidal
\(\pi\in\mathcal{V}^{\mathrm{prim}}\) (Hoffstein--Lockhart 1994
Ann.~Math.~140; Shahidi 1990 Ann.~Math.~132).
\item The adelic product
\(\prod_{p}L_{p}(1,\mathrm{Ad})=L(1,\mathrm{Ad})\) converges
absolutely on the cuspidal sph-fin locus (Langlands--Shahidi
method; Shahidi 1981 Amer.~J.~Math.~103).
\item The Koszul--Petersson compatibility
(\Cref{v4-arith:kp:thm:main-full}) holds on the normalized sph-fin
cuspidal core under the Koszul-dual normalisation, conditional on
\Cref{v4-arith:bcglue-k5:conj:local-comparison-datum};
full-sector extension (including Eisenstein handling beyond the
sph-fin cut) remains open.
\end{itemize}
\end{theorem}

\begin{proof}
\emph{Convergence of \(L(1,\mathrm{Ad})\).} Shahidi 1981 Amer.~J.~Math.~103
(``On certain \(L\)-functions'') proves that \(L(s,\mathrm{Ad})\) for
cuspidal \(\pi\) on \(GL_{n}/\mathbb{Q}\) has holomorphic continuation
to \(\mathrm{Re}(s)>1/2\) and satisfies a Langlands--Shahidi
functional equation. At \(s=1\), \(L(1,\mathrm{Ad})\) is finite;
Hoffstein--Lockhart 1994 Ann.~Math.~140 proves nonvanishing at
\(s=1\) (the adjoint \(L\)-function is nonvanishing on the edge of
the critical strip). Hence \(L(1,\mathrm{Ad})\in\mathbb{C}^{\times}\).

\emph{Adelic product convergence.} The Euler product
\(L(s,\mathrm{Ad})=\prod_{p}L_{p}(s,\mathrm{Ad})\) converges for
\(\mathrm{Re}(s)>1\) absolutely (standard Euler product argument;
each \(L_{p}(s,\mathrm{Ad})\) is a product of three geometric series
\((1-\alpha_{p}\beta_{p}^{-1}p^{-s})^{-1}\cdots\)). At \(s=1\),
\(L(1,\mathrm{Ad})\) is understood via analytic continuation from
\(\mathrm{Re}(s)>1\), with the continuation supplied by
Langlands--Shahidi. The adelic product equals the continued
value, a finite nonzero constant.

\emph{Rescaling preserves Koszul--Petersson compatibility.} The
Koszul--Petersson identity
(\Cref{v4-arith:kp:thm:main-full}, (\ref{v4-arith:kp:eq:main}))
reads
\[
\langle v,w\rangle^{\mathrm{Pet}}
\;=\;
\langle v,\sigma^{\mathrm{BC}}(w)\rangle^{\mathrm{Koszul}}.
\]
Under the rescaling
\(\langle\,\cdot\,,\,\cdot\,\rangle^{\mathrm{Koszul,Ad}}=L(1,\mathrm{Ad})^{-1}\cdot\langle\,\cdot\,,\,\cdot\,\rangle^{\mathrm{Koszul}}\),
the adjoint factor is absorbed on the Koszul side.  The corresponding
Petersson reduced pairing is a normalisation datum, not a consequence
of Bump \S3.8 alone.  Under this datum, the identity
\[
\langle v,w\rangle^{\mathrm{Pet,reduced}}
\;=\;
\langle v,\sigma^{\mathrm{BC}}(w)\rangle^{\mathrm{Koszul,Ad}}
\]
holds on the sph-fin cuspidal locus.

\emph{Tamagawa consistency (K6 comparison).} The Weil-normalised
Haar on \(GL_{2}(\mathbb{A})\) has Tamagawa volume
\(\mathrm{vol}([GL_{2}])=1\) (Weil 1967 Ch.~VII \S3). The product
of local volumes
\(\prod_{p}\mathrm{vol}(GL_{2}(\mathbb{Z}_{p}))\cdot\mathrm{vol}(GL_{2}(\mathbb{R})/GL_{2}(\mathbb{Z}))\)
equals \(1\) by Weil's Tamagawa computation; the Koszul-dual
normalisation by \(L(1,\mathrm{Ad})^{-1}\) is compatible with this
Tamagawa normalisation because \(L(1,\mathrm{Ad})\) is constant (a
single scalar across all primes) and the local volumes already
account for the \(\zeta_{p}(1)\zeta_{p}(2)^{-1}\) factor. Hence the
Koszul-dual-normalised pairing agrees with the Petersson pairing
after removing the \(L(1,\mathrm{Ad})\) factor, adelically
consistently.

K5 closes on the cuspidal sph-fin locus under the explicit convention of
\Cref{v4-arith:bcglue-k5:def:koszul-dual-norm}.
\end{proof}

\begin{remark}[K5 form upgrade]
\label{v4-arith:bcglue-k5:rem:K5-delta}
\Cref{v4-arith:sw:thm:KP-split}(K5) stated K5 as
``Theorem conditional on Godement--Jacquet absorption convention.''
\Cref{v4-arith:bcglue-k5:def:koszul-dual-norm} specifies the
convention explicitly: rescale the Koszul pairing by
\(L(s,\mathrm{Ad})^{-1}\) evaluated at \(s=1\).
With this specification,
\Cref{v4-arith:bcglue-k5:thm:K5-absorption-canonical} and
\Cref{v4-arith:bcglue-k5:thm:K5-adelic-consistent}
reduce K5 on the cuspidal sph-fin locus to the local comparison
datum. The
Hoffstein--Lockhart nonvanishing
\(L(1,\mathrm{Ad})\neq 0\) makes the convention well-defined on
the cuspidal sph-fin locus.
\end{remark}

\section{Consequence: K-P Compat, Conditional Domain}
\label{v4-arith:bcglue-k5:K-P-unconditional}

\begin{corollary}[K-P compat on the cuspidal sph-fin locus under the local comparison datum]
\label{v4-arith:bcglue-k5:cor:K-P-unconditional}
\ClaimStatusProvedHereConditional \emph{(conditional on
\Cref{v4-arith:bcglue-k5:conj:local-comparison-datum} and
domain-restricted to the cuspidal sph-fin locus)}. The Koszul--Petersson compatibility
(\Cref{v4-arith:kp:thm:main-full}) holds on the cuspidal sph-fin
locus of \(\mathcal{V}^{\mathrm{prim}}\), with the Koszul pairing
interpreted under the Koszul-dual normalisation convention of
\Cref{v4-arith:bcglue-k5:def:koszul-dual-norm}. Full-sector
extension further separates into the KMS\(_{1}\)/GNS half-density
statement of \Cref{v4-arith:kms-gns:thm:KMS1-GNS-BC-identification}
and the ramified Rankin--Selberg and Eisenstein conditions of
\Cref{v4-arith:kp:rem:domain-restriction}.
\end{corollary}

\begin{proof}
\Cref{v4-arith:kp:thm:main-full} reduced K-P compat to the seven
sub-claims K1--K7 of \Cref{v4-arith:sw:thm:KP-split}; K1, K3 are
definitional, K2, K4, K6, K7 are theorems by primary literature,
K5 was conditional on the GJ absorption convention. The preceding
calculation reduces K5 on the cuspidal sph-fin locus to
\Cref{v4-arith:bcglue-k5:thm:K5-absorption-canonical} and
\Cref{v4-arith:bcglue-k5:thm:K5-adelic-consistent}. All seven
sub-claims close on the cuspidal sph-fin locus after the local
comparison datum is assumed. Hence K-P compat closes on that
conditional domain; full-sector extension is the open residual documented in
\Cref{v4-arith:kp:rem:domain-restriction}.
\end{proof}

\begin{remark}[impact on P3]
\label{v4-arith:bcglue-k5:rem:P3-impact}
\Cref{v4-arith:kp:cor:P3-unconditional} stated P3 conditional on
K-P compat, which was itself conditional on K5
(\Cref{v4-arith:sw:thm:KP-split}).
\Cref{v4-arith:bcglue-k5:thm:K5-absorption-canonical} reduces K5 on
the cuspidal sph-fin locus to the local comparison datum. The logical
chain
\[
\text{K5 (cuspidal sph-fin)}
\;\Rightarrow\;
\text{K-P compat (cuspidal sph-fin)}
\;\Rightarrow\;
\text{P3 (cuspidal sph-fin)}
\]
is complete on the cuspidal sph-fin locus with the Koszul-dual
normalisation convention explicit.
\end{remark}

\begin{remark}[impact on sufficient RH path]
\label{v4-arith:bcglue-k5:rem:RH-path-impact}
\Cref{v4-arith:kp:cor:rh-minimal-updated} reduced the sufficient RH
path from three residuals (arith.~Pos., K-P compat, VB\(^{*}\)) to
two (arith.~Pos., VB\(^{*}\)) after K-P compat was established.
The K5 closure does not remove additional residuals
from the RH path but upgrades the form of K-P compat from
``conditional on convention'' to the explicit domain restriction of
\Cref{v4-arith:bcglue-k5:cor:K-P-unconditional}. The two remaining
residuals on the RH path, arith.~Positselski and VB\(^{*}\), are
unaffected and remain open.
\end{remark}

\section{Residual Conditions}
\label{v4-arith:bcglue-k5:residual}

The open residuals split as follows.

\begin{description}
\item[(BCglue.local).] Conditional;
\Cref{v4-arith:bcglue-k5:thm:local-BZCHN-plus-EG}.
Residual condition: non-formal theorem.
\item[(BCglue.glue).] Conjectural (prismatic flatness);
\Cref{v4-arith:bcglue-k5:thm:prismatic-gluing}.
Residual condition: non-formal theorem.
\item[(BC-diamond-glue.core).] Residual conjecture conditional on
arithmetic AG (\Cref{v4-arith:thm:arithmetic-AG}). Domain:
non-formal theorem past arithmetic AG.
\item[Arithmetic AG (global).] Open conjecture. Domain: substantial
(primary-literature-length); this is
the cumulative domain of extending Fargues--Scholze local
geometrization to a global arithmetic statement, estimated
non-formal theorem past Fargues--Scholze.
\item[K5 (GJ absorption convention).] Conditional on
\Cref{v4-arith:bcglue-k5:conj:local-comparison-datum} on the cuspidal
sph-fin locus by
\Cref{v4-arith:bcglue-k5:thm:K5-absorption-canonical} and
\Cref{v4-arith:bcglue-k5:thm:K5-adelic-consistent}. Remaining condition:
the local comparison datum. Eisenstein handling via the
sph-fin cut of \Cref{v4-arith:sw:thm:KP-split}(K7).
\item[K-P compat, P3.] Domain-restricted to the cuspidal sph-fin
locus by \Cref{v4-arith:bcglue-k5:cor:K-P-unconditional}. Domain
conditional on the local comparison datum. The full-sector extension
has the KMS\(_{1}\)/GNS half-density factor separated from the
ramified Rankin--Selberg and Eisenstein conditions.
\item[Sufficient RH path.] Four-input conditional conjunction
(\Cref{v4-arith:w3-comparison-B:thm:four-input-branch}): arith.~Positselski
(five open sub-items), Koszul--Pet compat (domain-restricted),
H-P\(^{\mathrm{Ar}}\), VB\(^{*}\) main direction.
Arith.~Positselski: \(\sim 30\)--\(50\) hypothesis.
VB\(^{*}\) is the classical positive-trace open problem.
\end{description}

\begin{remark}[residual form]
\label{v4-arith:bcglue-k5:rem:residual-summary}
The residual form is:
\begin{enumerate}
\item (BCglue.local): non-formal theorem, form \emph{conditional}; the
integration of Ben-Zvi--Chen--Helm--Nadler + Emerton--Gee + Zhu +
Fargues--Scholze into a single canonical local coherent Springer sheaf,
including \(\ell=p\) compatibility and \(\ell\)-independence, remains a
primary-source writing condition; see
\Cref{v4-arith:bcglue-k5:thm:local-BZCHN-plus-EG}.
\item (BCglue.glue): non-formal theorem, form \emph{conjectural} conditional
on the prismatic flatness of
\Cref{v4-arith:bcglue-k5:lem:prismatic-flatness}; the cross-prime
Nygaard-convention-aware transition theorem is not in the primary
literature.
\item K5 (GJ absorption): reduced on the cuspidal sph-fin locus to
the local comparison datum under the Koszul-dual normalisation convention
(\Cref{v4-arith:bcglue-k5:def:koszul-dual-norm}); Eisenstein
handling via the sph-fin cut.
\item P3 form restriction: closed, domain-restricted theorem on
the cuspidal sph-fin locus with the explicit normalisation
convention; full-sector P3 remains open.
\item (BC-diamond-glue): a single cross-prime compatibility remains,
non-formal theorem past arithmetic AG.
\item Arithmetic AG, arithmetic Positselski, and VB\(^{*}\) are not
logical consequences of K5.
\end{enumerate}
\end{remark}

\section{Residual Boundary}
\label{v4-arith:bcglue-k5:summary}

\begin{enumerate}
\item (BC-diamond-glue) decomposes into three sub-operations:
(BCglue.local), (BCglue.glue), (BC-diamond-glue.core)
(\Cref{v4-arith:bcglue-k5:thm:decomposition}). The decomposition
is logically independent.
\item (BCglue.local) is conditional: Ben-Zvi--Chen--Helm--Nadler +
Emerton--Gee + Zhu + Fargues--Scholze supply the inputs, but a single
canonical primary-source transition theorem covering \(\ell = p\) and
\(\ell\)-independence has not yet been written
(\Cref{v4-arith:bcglue-k5:thm:local-BZCHN-plus-EG}).
\item (BCglue.glue) is conjectural conditional on prismatic flatness
of the Nygaard-convention-aware transition maps
(Bhatt--Scholze, Bhatt--Lurie; see
\Cref{v4-arith:bcglue-k5:lem:prismatic-flatness}) and filtered
colimits (Gaitsgory--Rozenblyum)
(\Cref{v4-arith:bcglue-k5:thm:prismatic-gluing}).
\item The irreducible residual of (BC-diamond-glue) is
(BC-diamond-glue.core): a non-formal theorem cross-prime compatibility,
conditional on arithmetic AG
(\Cref{v4-arith:bcglue-k5:thm:core-compatibility}).
\item \(\mathrm{L}_{\mathrm{ACD}}\) tight residual list
(\Cref{v4-arith:bcglue-k5:cor:LACD-tight-list}): (BCglue.local)
conditional; (BCglue.glue) conjectural; (BC-diamond-glue.core)
conjectural conditional on arithmetic AG. The three subconditions
together form the residual content of \(\mathrm{L}_{\mathrm{ACD}}\)
past F2.c and the subconditions of
Ch.~\ref{v4-arith:lacd-sub}.
\item K5 is reduced on the cuspidal sph-fin locus to the local
comparison datum and the Koszul-dual normalisation convention
(\Cref{v4-arith:bcglue-k5:def:koszul-dual-norm}): rescale the
Koszul pairing at each finite prime by
\(L_{p}(1,\mathrm{Ad})^{-1}\), giving a global rescaling by
\(L(1,\mathrm{Ad})^{-1}\). The rescaling is canonical on the
cuspidal sph-fin locus, and Hoffstein--Lockhart 1994 + Shahidi
1990 guarantee \(L(1,\mathrm{Ad})\neq 0\) on the cuspidal sph-fin
locus
(\Cref{v4-arith:bcglue-k5:thm:K5-absorption-canonical},
\Cref{v4-arith:bcglue-k5:thm:K5-adelic-consistent}); Eisenstein
handling is via the sph-fin cut of
\Cref{v4-arith:sw:thm:KP-split}(K7).
\item K-P compat and P3 are conditional domain-restricted statements
on the cuspidal sph-fin locus with the explicit normalisation convention
(\Cref{v4-arith:bcglue-k5:cor:K-P-unconditional});
the full-sector extension separates the KMS\(_{1}\)/GNS half-density
from the ramified Rankin--Selberg and Eisenstein conditions.
\item Residual condition: (BC-diamond-glue) has one irreducible
cross-prime compatibility past arithmetic AG; the primary open
residuals remain arithmetic AG, arith.~Positselski (five open
sub-items), H-P\(^{\mathrm{Ar}}\), and VB\(^{*}\) (the
four-input conditional conjunction of
\Cref{v4-arith:w3-comparison-B:thm:four-input-branch}).
\end{enumerate}

(BC-diamond-glue) is isolated as a single cross-prime compatibility
conditional on arithmetic AG, and K5 is reduced on the cuspidal
sph-fin locus to the local comparison datum. The K-P compatibility
and P3 are conditional domain-restricted statements on the cuspidal
sph-fin locus. The sufficient RH path is unaffected.

The primary sources are drawn from seven disjoint lists:
Ben-Zvi--Chen--Helm--Nadler (coherent Springer theory);
Emerton--Gee (Galois deformation stacks); Bhatt--Scholze +
Bhatt--Lurie (prismatic cohomology and Nygaard convention);
Gaitsgory--Rozenblyum (derived algebraic geometry);
Godement--Jacquet (local functional equation); Langlands--Shahidi
(Shahidi 1990, 1981; adjoint \(L\)-functions and nonvanishing);
Bost--Connes (KMS\(_{1}\) modular theory with the Tomita
half-density).
No pair of lists shares a primary author at overlapping work;
the independence of sources requirement of the comparison
is satisfied.
